\textbf{Proof.}
By \(\mathrm{def:bias\text{-}aware\text{-}interval}\), the event
\(\tau\in\mathcal I_{h,J_{\mu,n},J_{r,n}}\) is exactly
\[
\left|\widehat\tau^{\mathrm{sp}}_{h,J_{\mu,n},J_{r,n}}-\tau\right|
\le\overline B_X(h,J_{\mu,n},J_{r,n})
+R_S(h,J_{\mu,n},J_{r,n})+R_T(m).
\tag{1}
\]
For each \(\mathbb P_{n,m}\in\mathcal M_{n,m}(a,s,t)\),
\(\mathrm{thm:ricot\text{-}attainment}\) assigns probability at least
\(1-\gamma\) to (1).  The constants and the H\"older class in this assertion are
fixed uniformly across the displayed model class.  Taking the infimum over its
members therefore gives
\[
\inf_{\mathbb P_{n,m}\in\mathcal M_{n,m}(a,s,t)}
\mathbb P_{n,m}\{\tau\in\mathcal I_{h,J_{\mu,n},J_{r,n}}\}
\ge1-\gamma.
\tag{2}
\]

The interval is symmetric about
\(\widehat\tau^{\mathrm{sp}}_{h,J_{\mu,n},J_{r,n}}\), and its half-length is
\(\overline B_X+R_S+R_T\).  Its realized length is consequently exactly
\[
2\{\overline B_X(h,J_{\mu,n},J_{r,n})
+R_S(h,J_{\mu,n},J_{r,n})+R_T(m)\}.
\tag{3}
\]
If the realized covariate design satisfies \(\overline B_X\le b\) and
\(\sum_i\Gamma_i^2\le g^2\), then
\(\mathrm{def:foldwise\text{-}inference\text{-}statistics}\), as incorporated in
the interval construction, gives
\[
R_S\le\varepsilon^{-1}\sqrt{2\log(4/\gamma)}\,g,
\qquad
R_T(m)=\sqrt{2\log(4/\gamma)/m}.
\tag{4}
\]
Substitution of (4) into (3) proves the claimed realized-design length bound.

Finally, \(\overline B_X\) is, by the construction in
\(\mathrm{def:bias\text{-}envelope}\), the maximum of two finite signed jet linear
programs.  If \(N\) is the number of distinct observed covariate sites, each program
has one variable \(u_{\ell,r}\) for every
\(1\le\ell\le N\) and \(0\le r\le p\), hence \(N(p+1)\) variables.  The pairwise
jet restrictions contribute \(O\{N^2(p+1)\}\) one-sided linear inequalities after
each absolute-value constraint is split into its two signs; the box restrictions
do not change this order.  This verifies the stated finite computation.  No
averaging or stochastic-order control of \(\overline B_X\) or
\(\sum_i\Gamma_i^2\) was used, so the proof entails none of the expected-length,
Gaussian-approximation, multiplier-approximation, variance-consistency, leverage,
or frontier-attainment claims expressly excluded by the theorem.
\(\square\)